۱۴۴  
من در گردبادی از یادگیری غرق شده بودم. با دوستانم تماس نگرفتم. به‌ندرت با خانواده‌ام صحبت می‌کردم. هر پیشنهادی برای نوشتن که به من می‌شد را رد کردم. من در یک واقعیت دیگر زندگی می‌کردم.  
«به راسپوتین گفتم»، استیو گفت یک شب، «که بیش از همه‌ی پسرهای اغواگر دیگر، دوست داشتم تو یکی از مربیان ما شوی.»

پیشنهادی بود که باید ردش می‌کردم. دنیای اغواگری قصری بود از درهای باز. قدم گذاشتن از یکی از آن‌ها، بدون توجه به گنج‌های وسوسه‌انگیز درونش، به معنی بستن بقیه‌ی درها بود.  
یک شنبه شب از سن دیگو به خانه برگشتم و پیغامی از کلیف، از لیست کلیف، روی دستگاه بام پیدا کردم. او در شهر بود و می‌خواست مرا با آخرین کشف PUA‌اش آشنا کند — موتورسوار سابقی که کارگر ساختمانی شده بود و خودش را دیوید X می‌نامید.  
کلیف از ابتدای تشکیل جامعه‌اش در آنجا بود. او در چهل سالگی بود و به اندازه‌ی سختگیر بودنش خوش‌برخورد بود. با اینکه ظاهر جذاب سنتی داشت، ولی تجسم زنده‌ی کلمه‌ی «چهارخانه» بود. او شبیه کسی بود که از یک سریال خانوادگی دهه ۵۰ بیرون آمده باشد. او ادعا می‌کرد که در خانه‌اش یک کمد با بیش از هزار کتاب اغواگری دارد. نسخه‌های مجله Pick-Up Times، یا همان مجله کوتاه‌مدت دهه ۷۰ و یک نسخه اصیل از کتاب کلاسیک اریک وبر با عنوان «چگونه دخترها را اغوا کنیم» و کتاب‌های کج‌معنا با عناوین مبهم مانند «اغواگری زمانی شروع می‌شود که زن می‌گوید نه» در آنجا بود.  
دیوید X یکی از نیم‌دوجین PUA‌هایی بود که کلیف در طول سال‌ها کشف کرده و در لیست خود تبلیغ کرده بود، لیستی که او بعد از اینکه راس در لیست پستی Speed Seduction او را به خاطر بحث درباره‌ٔ تکنیک اغواگری که به NLP مرتبط نبود، مورد انتقاد قرار داده بود، در سال ۱۹۹۹ شروع کرد. هر PUA یک تخصص داشت و تخصص دیوید X مدیریت حرم‌سرا بود — مدیریت روابط با چندین زن بدون دروغ گفتن به آن‌ها.  
وقتی وارد رستوران دیم‌سام شدیم، از آنچه که دیدم شوکه شدم. دیوید X شاید زشت‌ترین PUA‌ای بود که تا به حال ملاقات کرده بودم. او به‌طرز وحشتناکی بزرگ، کچل و شبیه قورباغه بود، با زگیل‌هایی که صورتش را پوشانده بودند و صدایی شبیه به صد هزار پک سیگار.  
غذایم با او شبیه بسیاری از غذاهایی بود که قبلاً داشتم. جز اینکه قوانین همیشه متفاوت بود. قوانین او چنین بود:  
۱. چه اهمیتی دارد او چه فکر می‌کند؟

۲. تو مهم‌ترین شخص در این رابطه هستی.  

فلسفه او این بود که هرگز به یک زن دروغ نگوید. او به خودش می‌بالید که با به دام انداختن زنان با کلمات خودشان با آن‌ها هم‌خوابه می‌شود. به‌عنوان مثال، در ملاقات با دختری در یک بار، او را وادار می‌کرد که بگوید خودجوش است و هیچ‌قانونی ندارد؛ سپس، اگر او تمایلی به ترک بار با او نداشت، می‌گفت: «فکر کردم تو خودجوش هستی. فکر کردم هر کاری که بخواهی انجام می‌دهی.»  
او روی صندلی‌اش به مانند تکه‌ای از پنیر سوئیسی در حال ذوب شدن پهن شد و به ما اطلاع داد: «تنها دروغ‌هایی که می‌گویم اینست که: 'در دهانت نمی‌آیم' و 'فقط دور باسنت می‌مالم'.» این یک تصویر زیبا نبود.  

فلسفه او با تمامی آنچه از مستری آموخته بودم، در تضاد کامل بود، و او در طول شام به من این را فهماند. او شاهدی بر نظریه‌ی دهان بزرگ کلیف، یک مرد آلفای طبیعی و سخنگو بود.

«بهترین چیز این است»، او با افتخار گفت، «افرادی شبیه به من و افرادی شبیه به تو و مستری در جهان هستند. وقتی شما هنوز در بار در حال انجام حقه‌های جادویی هستید، من برای دور دوم برگشته‌ام.»

یک شام جالب بود و من چیزهای کوچک بسیاری درباره‌ی بازی آموختم که بارها از آن‌ها استفاده می‌کردم. اما تا زمان پایان صبحانه، متوجه چیزی شدم: نیازی به ملاقات با هر گونه استاد دیگری ندارم.

من هر اطلاعاتی را که برای تبدیل شدن به بزرگترین هنرمند اغوا در جهان نیاز داشتم، داشتم. صدها گشایش‌کننده، روش‌، جملات شوخ و با اعتماد‌به‌نفس، روش‌های نشان دادن ارزش و تکنیک‌های جنسی قوی داشتم. و به والهالا و دوباره هیپنوتیزم شده بودم. نیازی به یادگیری چیز دیگری نبود، مگر برای لذت و علاقه شخصی خودم. فقط باید در میدان بودم - مدام نزدیک می‌شدم، کالیبره می‌کردم، تنظیمات نهایی انجام می‌دادم، و از موانع می‌گذشتم. آمادهٔ میامی بودم و همهٔ کارگاه‌های بعدی.

وقتی کلیف مرا به خانه رساند، به خودم وعده‌ای دادم: اگر بار دیگر با استادی ملاقات کنم، این ملاقات باید به‌عنوان همتا باشد نه به‌عنوان شاگرد.